\chapter{Sectional Curvature Comparison II}
\label{chap:pet-sectional-curvature-comparison-ii}

Faithful transcription of Petersen, \emph{Riemannian Geometry} (GTM 171, 3rd ed.),
Chapter 12. The chapter develops the global consequences of a lower sectional
curvature bound through the critical point theory of distance functions
initiated by Grove and Shiohama (§12.1) and Toponogov's comparison theorem
(§12.2), with applications to the quarter-pinched and diameter sphere theorems
(§12.3), the Cheeger--Gromoll/Gromoll--Meyer soul theorem for open manifolds of
nonnegative curvature (§12.4), Gromov's finiteness of Betti numbers (§12.5),
and the Grove--Petersen homotopy finiteness theorem (§12.6), together with a
guide to further reading (§12.7) and the chapter's exercises (§12.8). Each
numbered statement keeps Petersen's lemma/proposition/theorem number in the
environment title; proofs keep his explicit cross-references to earlier
chapters so the dependency graph is recoverable from the text.

\section{Critical Point Theory}

The idea of triangle comparison for surfaces goes back to Alexandrov, who
influenced Toponogov, although Pizzetti had already established local triangle
comparison on surfaces at the start of the 20th century. In the generalized
critical point theory of this section, the object is to define generalized
gradients of continuous functions and use them to show that certain regions of
a manifold carry no topology. The motivating case is the following classical
fact about smooth proper functions.

\begin{lemma}[Lem. 12.1.1 — deformation retraction along a gradient flow]
  \label{lem:pet-ch12-smooth-deformation-retraction}
  \lean{PetersenLib.gradientFlowDeformationRetraction}
  Let $(M,g)$ be a Riemannian manifold and $f:M\to\mathbb{R}$ a proper smooth
  function. If $f$ has no critical values in the closed interval $[a,b]$, then
  $f^{-1}((-\infty,b])$ and $f^{-1}((-\infty,a])$ are diffeomorphic; more
  precisely, there is a deformation retraction of $f^{-1}((-\infty,b])$ onto
  $f^{-1}((-\infty,a])$, so the inclusion
  $f^{-1}((-\infty,a])\hookrightarrow f^{-1}((-\infty,b])$ is a homotopy
  equivalence.
  \source{petersen:page-0457,petersen:page-0458}
\end{lemma}
\begin{proof}
  \uses{lem:pet-ch12-smooth-deformation-retraction}
  Since $f$ has no critical points on the compact set $f^{-1}([a,b])$,
  $\nabla f$ is nonzero there. Let $\psi:M\to[0,1]$ be a bump function equal to
  $1$ on $f^{-1}([a,b])$ and $0$ outside a compact neighborhood of it, and set
  \[
    X=-\psi\cdot\frac{\nabla f}{|\nabla f|^2}.
  \]
  Since $X$ has compact support it is complete; let $F^t$ be its flow. For
  fixed $q\in M$, $\tfrac{d}{dt}f(F^t(q))=g(X,\nabla f)$, which equals $-1$
  wherever $\psi=1$, so as long as the integral curve $t\mapsto F^t(q)$ remains
  in $f^{-1}([a,b])$, $t\mapsto f(F^t(q))$ is affine with slope $-1$.
  Consequently $F^{b-a}:M\to M$ carries $f^{-1}((-\infty,b])$
  diffeomorphically into $f^{-1}((-\infty,a])$, and
  \[
    r_t:f^{-1}((-\infty,b])\to f^{-1}((-\infty,b]),\qquad
    r_t(p)=\begin{cases}p&f(p)\le a\\F^{t(f(p)-a)}(p)&a\le f(p)\le b\end{cases}
  \]
  is a well-defined homotopy with $r_0=\mathrm{id}$ and $r_1$ a retraction of
  $f^{-1}((-\infty,b])$ onto $f^{-1}((-\infty,a])$, exhibiting the inclusion of
  the latter into the former as a homotopy equivalence. (Properness of $f$ is
  used essentially to make $X$ complete: deleting a single point from
  $f^{-1}([a,b])$ leaves $f$ still without critical values there, yet the
  conclusion visibly fails.)
\end{proof}

We generalize \cref{lem:pet-ch12-smooth-deformation-retraction} to distance
functions, which need not be $C^1$. Suppose $(M,g)$ is complete and
$K\subset M$ is compact; then $r(x)=|xK|=\min\{|xp|\mid p\in K\}$ is proper.
Wherever $r$ is smooth it has unit gradient and hence is noncritical there,
but $r$ may also have local maxima, where one still wants a notion of
noncriticality. This is provided by the following generalized gradient.

\begin{definition}[segment direction sets; regular, critical, and
  $\alpha$-regular points]
  \label{def:pet-ch12-alpha-regular}
  \lean{PetersenLib.SegmentDirectionSet,PetersenLib.IsAlphaRegularAt,PetersenLib.IsRegularAt}
  Let $(M,g)$ be complete, $K\subset M$ compact, and $r(x)=|xK|$. For
  $x\in M$ let $\overrightarrow{xK}\subset T_xM$ denote the set of initial
  unit velocities of unit-speed minimizing segments from $x$ to $K$ (i.e.\ of
  length $r(x)$); when $K=\{q\}$ this is written $\overrightarrow{xq}$. If $r$
  is smooth at $x$ then $\overrightarrow{xK}$ has a single element and
  $-\nabla r(x)=\overrightarrow{xK}$. In general $x$ is \emph{regular}, or
  \emph{noncritical}, for $r$ if $\overrightarrow{xK}$ lies in an open
  hemisphere of the unit sphere of $T_xM$, equivalently there is a unit
  $v\in T_xM$ with $\angle(v,\overrightarrow{xK})<\pi/2$ for every
  $\overrightarrow{xK}\in\overrightarrow{xK}$ (any such $v$ is the center of
  such a hemisphere); otherwise $x$ is a \emph{critical point} of $r$. More
  generally, for $\alpha\in(0,\pi]$, $x$ is \emph{$\alpha$-regular} for $r$ if
  there is a unit $v\in T_xM$ with $\angle(v,\overrightarrow{xK})<\alpha$ for
  every $\overrightarrow{xK}\in\overrightarrow{xK}$ (so regular means
  $\pi/2$-regular), and $R_\alpha(x,K)\subset T_xM$ is the set of all such
  unit directions $v$ at an $\alpha$-regular point $x$.
  \source{petersen:page-0458}
\end{definition}

It was Berger who first showed that a local maximum of $r$ must be critical
in this sense; his result is a consequence of the following proposition.

\begin{proposition}[Prop. 12.1.2 — properties of regular points]
  \label{prop:pet-ch12-regular-point-properties}
  \lean{PetersenLib.regularPointProperties}
  \uses{def:pet-ch12-alpha-regular}
  Suppose $(M,g)$ is complete, $K\subset M$ compact, $r(x)=|xK|$. Then:
  \begin{enumerate}
    \item $\overrightarrow{xK}$ is closed, hence compact, for every $x\in M$.
    \item The set of $\alpha$-regular points is open in $M$.
    \item $R_\alpha(x,K)$ is convex for every $\alpha\le\pi/2$.
    \item If $U\subset M$ is open and every point of $U$ is $\alpha$-regular,
      there is a unit vector field $X$ on $U$ with $X|_x\in R_\alpha(x,K)$ for
      all $x\in U$; moreover, if $c$ is an integral curve of $X$ and $s<t$,
      \[
        |c(s)K|-|c(t)K|\ge\cos(\alpha)(t-s).
      \]
  \end{enumerate}
  \source{petersen:page-0459,petersen:page-0460}
\end{proposition}
\begin{proof}
  \uses{prop:pet-ch12-regular-point-properties}
  \emph{(1)} If $\overrightarrow{xK_i}\in\overrightarrow{xK}$ is a sequence of
  initial velocities of minimizing segments from $x$ to $K$ converging to a
  unit $v\in T_xM$, then $\exp_x(tv)$, $t\in[0,r(x)]$, is the limit of these
  segments and hence itself a segment realizing $|xK|$; since $K$ is closed,
  $\exp_x(r(x)v)\in K$, so $v\in\overrightarrow{xK}$. Thus
  $\overrightarrow{xK}$ is closed, hence compact as a closed subset of the
  unit sphere in $T_xM$.

  \emph{(2)} Suppose $x_i\to x$ with each $x_i$ not $\alpha$-regular; we show
  $x$ is not $\alpha$-regular. Fix a unit $v\in T_xM$ and choose $v_i\in
  T_{x_i}M$ with $v_i\to v$. By hypothesis there is, for each $i$, some
  $\overrightarrow{x_iK}\in\overrightarrow{x_iK}$ with
  $\angle(v_i,\overrightarrow{x_iK})\ge\alpha$; passing to a subsequence these
  converge to a unit $w\in T_xM$ with $\angle(v,w)\ge\alpha$. The segments
  $t\mapsto\exp_{x_i}(t\,\overrightarrow{x_iK})$, $t\in[0,|x_iK|]$, then
  converge to $\exp_x(tw)$, $t\in[0,|xK|]$, forcing the latter to be a segment
  from $x$ to $K$, i.e.\ $w\in\overrightarrow{xK}$. As $v$ was an arbitrary
  unit vector, no $v$ works, so $x$ is not $\alpha$-regular.

  \emph{(3)} For $w\in T_xM$ the open cone
  $C_\alpha(w)=\{v\in T_xM\mid\angle(v,w)<\alpha\}$ is convex whenever
  $\alpha\le\pi/2$ (an elementary property of the round unit sphere). Since
  $R_\alpha(x,K)=\bigcap_{\overrightarrow{xK}\in\overrightarrow{xK}}
  C_\alpha(\overrightarrow{xK})$ is an intersection of such cones, it is
  convex.

  \emph{(4)} For $p\in U$ choose $v_p\in R_\alpha(p,K)$ and extend it to a
  unit vector field $V_p$; by the proof of (2), $V_p(x)\in R_\alpha(x,K)$ for
  $x$ near $p$, so $V_p$ is a generalized gradient direction on a neighborhood
  $U_p$ of $p$. Choose a locally finite subcover $\{U_i\}$ of $U$ by such
  neighborhoods with subordinate partition of unity $\{\lambda_i\}$; by
  convexity (part (3)), $V=\sum_i\lambda_iV_i$ is nonzero on $U$ and
  $X=V/|V|$ still satisfies $X|_x\in R_\alpha(x,K)$ for every $x\in U$.

  For the length estimate it is convenient to work with $-r$. Where $r$ is
  smooth at $c(s)=x$, $\tfrac{d}{ds}(-r\circ c)(s)=g(X,-\nabla r)
  =\cos\angle(X,\overrightarrow{xK})>\cos\alpha$, directly from
  $X|_x\in R_\alpha(x,K)$. In general $-r\circ c$ is Lipschitz, hence
  absolutely continuous and almost-everywhere differentiable. At a point
  $x=c(0)$ where $\nabla r$ may fail to exist, choose a variation
  $\tilde c(s,t)$ with $t\mapsto\tilde c(0,t)$ a unit-speed segment from $K$
  to $x$, $\tilde c(s,0)=\tilde c(0,0)\in K$, $|\partial_t\tilde c(s,t)|$
  independent of $t$ (equal to the $t$-curve's length), and
  $\tilde c(s,1)=c(s)$. With $E(\tilde c_s)=\tfrac12\int_0^1|\partial_t\tilde
  c|^2\,dt$ the energy, Cauchy--Schwarz gives
  \[
    \tfrac12\big(r(c(s))\big)^2
    \le\tfrac12\Big(\int_0^1|\partial_t\tilde c|\,dt\Big)^2
    \le E(\tilde c_s),
  \]
  with equality at $s=0$, so $E(\tilde c_s)$ is a support function from above
  for $\tfrac12(r\circ c)^2$ at $s=0$. If $r\circ c$ is differentiable at
  $s=0$, differentiating and applying the first-variation formula for energy
  along the geodesic $t\mapsto\tilde c(0,t)$ gives
  \[
    r(x)\,\tfrac{d(r\circ c)}{ds}\Big|_{s=0}=\tfrac{dE}{ds}\Big|_{s=0}
    =g\big(\partial_t\tilde c(0,1),\partial_s\tilde c(0,1)\big)
    =g\big(\partial_t\tilde c(0,1),X\big)
    =-r(x)\cos\angle\big(X,\overrightarrow{xK}\big),
  \]
  using $\partial_t\tilde c(0,1)=-r(x)\overrightarrow{xK}$ (the reversed
  terminal velocity of the segment from $K$ to $x$) for the last step; thus
  \[
    -\frac{d|c(s)K|}{ds}\Big|_{s=0}=\cos\angle\big(X,\overrightarrow{xK}\big)>\cos\alpha,
  \]
  the inequality holding since $X|_x\in R_\alpha(x,K)$. As this estimate holds
  at every point of differentiability of the absolutely continuous function
  $r\circ c$, integrating from $s$ to $t$ gives
  $|c(s)K|-|c(t)K|\ge\cos(\alpha)(t-s)$.
\end{proof}

We can now generalize \cref{lem:pet-ch12-smooth-deformation-retraction} to
distance functions.

\begin{lemma}[Lem. 12.1.3 — Grove--Shiohama retraction lemma]
  \label{lem:pet-ch12-retraction-lemma}
  \lean{PetersenLib.grovShiohamaRetraction}
  \uses{def:pet-ch12-alpha-regular,prop:pet-ch12-regular-point-properties}
  Let $(M,g)$ be complete and $r(x)=|xK|$ the distance to a compact submanifold
  $K\subset M$. If every point of $r^{-1}([a,b])$ is $\alpha$-regular for
  $\alpha<\pi/2$, then $r^{-1}((-\infty,b])$ deformation retracts onto
  $r^{-1}((-\infty,a])$.
  \source{petersen:page-0461}
\end{lemma}
\begin{proof}
  \uses{lem:pet-ch12-retraction-lemma,prop:pet-ch12-regular-point-properties}
  By \cref{prop:pet-ch12-regular-point-properties}(4) there is a compactly
  supported vector field $X$ on $r^{-1}([a,b])$ whose flow $F^t$ satisfies
  \[
    r(p)-r(F^t(p))>t\cos(\alpha),\qquad t\ge0,\ p,F^t(p)\in r^{-1}([a,b]).
  \]
  For each $p\in r^{-1}(b)$ there is a first time $t_p\le(b-a)/\cos\alpha$ with
  $F^{t_p}(p)\in r^{-1}(a)$, and $p\mapsto t_p$ is continuous. Hence
  \[
    r_t:r^{-1}((-\infty,b])\to r^{-1}((-\infty,b]),\qquad
    r_t(p)=\begin{cases}p & r(p)\le a\\ F^{t\cdot t_p}(p) & a\le r(p)\le b\end{cases}
  \]
  is a well-defined homotopy from the identity to a retraction of
  $r^{-1}((-\infty,b])$ onto $r^{-1}((-\infty,a])$.
\end{proof}

\begin{remark}[Rem. 12.1.4 — smooth approximation of the retraction]
  \label{rem:pet-ch12-retraction-remark}
  \lean{PetersenLib.grovShiohamaSmoothApproximation}
  \uses{lem:pet-ch12-retraction-lemma}
  The original Grove--Shiohama construction shows more: on the region where the
  distance function has no critical points it can be approximated by smooth
  functions without critical points, a fact that has been important elsewhere.
  \source{petersen:page-0461}
\end{remark}

\begin{corollary}[Cor. 12.1.5 — diffeomorphism to the normal bundle]
  \label{cor:pet-ch12-normal-bundle-diffeomorphism}
  \lean{PetersenLib.diffeomorphicToNormalBundle}
  \uses{lem:pet-ch12-retraction-lemma,def:pet-ch12-alpha-regular}
  Suppose $K$ is a compact submanifold of a complete Riemannian manifold $(M,g)$
  and $|xK|$ is regular everywhere on $M-K$. Then $M$ is diffeomorphic to the
  normal bundle of $K$ in $M$. In particular, if $K=\{p\}$, $M$ is diffeomorphic
  to $\mathbb{R}^n$.
  \source{petersen:page-0461,petersen:page-0462}
\end{corollary}
\begin{proof}
  \uses{cor:pet-ch12-normal-bundle-diffeomorphism}
  On $M-K$ there is a vector field $-X$ along which $|xK|$ increases, and near
  $K$ the distance function is smooth with $X=-x\vec K$. By the normal
  exponential map $\exp^+:T^+K\to M$ (a diffeomorphism from a neighborhood of the
  zero section onto a neighborhood of $K$, Cor.\ 5.5.3 of the source) and the
  fact that $t\mapsto\exp(tv)$ coincides for small $t$ with the integral curve of
  $-X$ starting in direction $v$, each $v\in T^+K$ determines a unique integral
  curve $c_v:(0,\infty)\to M$ of $-X$ with $\lim_{t\to0}\dot c_v(t)=v$. Define
  $F:T^+K\to M$ by $F(0_p)=p$ and $F(tv)=c_v(t)$ for $|v|=1$; this is
  differentiable and agrees with $\exp^+$ for small $t$. It is injective because
  integral curves of $-X$ cannot cross, and surjective because such integral
  curves must leave every sublevel of the proper function $|xK|$, hence are
  defined for all $t>0$. Finally $F$ is a diffeomorphism onto a neighborhood of
  $K$ by the normal exponential map, and elsewhere the flow of $-X$ acts by local
  diffeomorphisms, so $DF$ is everywhere nonsingular.
\end{proof}

\section{Distance Comparison}

Preparation for the Toponogov comparison theorem: given a Riemannian manifold
$(M,g)$, a \emph{hinge} consists of two segments $\overline{px}$, $\overline{xy}$
meeting at an interior angle $\alpha=\angle(\overrightarrow{xp},\overrightarrow{xy})$
at $x$, and a \emph{triangle} consists of three segments $\overline{xy}$,
$\overline{yz}$, $\overline{zx}$ meeting pairwise at $x,y,z$; geodesics may be
used in place of segments, allowing degenerate hinges or triangles with
coincident vertices.

\begin{definition}[hinge]
  \label{def:pet-ch12-hinge}
  \lean{PetersenLib.Hinge}
  A \emph{hinge} in $(M,g)$ consists of two segments (or geodesics)
  $\overline{px}$ and $\overline{xy}$ that form an interior angle
  $\alpha=\angle(\overrightarrow{xp},\overrightarrow{xy})$ at the common vertex
  $x$, where the specified directions are tangent to the given segments.
  \source{petersen:page-0462}
\end{definition}

\begin{definition}[triangle]
  \label{def:pet-ch12-triangle}
  \lean{PetersenLib.GeodesicTriangle}
  A \emph{triangle} in $(M,g)$ consists of three segments (or geodesics)
  $\overline{xy}$, $\overline{yz}$, $\overline{zx}$ meeting pairwise at the three
  vertices $x,y,z$.
  \source{petersen:page-0462,petersen:page-0463}
\end{definition}

\begin{lemma}[Lem. 12.2.1 — existence of comparison hinges/triangles]
  \label{lem:pet-ch12-comparison-hinge-triangle-exist}
  \lean{PetersenLib.existsComparisonHingeOrTriangle}
  \uses{def:pet-ch12-hinge,def:pet-ch12-triangle}
  Suppose $(M,g)$ is complete with $\sec\ge k$. Then for each hinge or triangle
  in $M$ there is a comparison hinge or triangle in the space form $S_k^n$ where
  the corresponding segments have the same lengths and, in the hinge case, the
  angle at the corresponding vertex is the same.
  \source{petersen:page-0463}
\end{lemma}
\begin{proof}
  \uses{lem:pet-ch12-comparison-hinge-triangle-exist}
  If $k>0$, $\mathrm{diam}\,M\le\pi/\sqrt k=\mathrm{diam}\,S_k^n$ (Cor.\ 6.3.2 of the
  source), so every segment involved has length $\le\pi/\sqrt k$.

  \emph{Hinge case.} Given segments $\overline{px}$, $\overline{xy}$ with interior
  angle $\alpha$ at $x$, choose a segment $\overline{p_kx_k}$ in $S_k^n$ of length
  $|px|$; at $x_k$ choose a direction $\overrightarrow{x_ky_k}$ with
  $\angle(\overrightarrow{x_kp_k},\overrightarrow{x_ky_k})=\alpha$, and along the
  unique geodesic in this direction select $y_k$ with $|x_ky_k|=|xy|$.

  \emph{Triangle case.} Choose $x_k,y_k$ with $|xy|=|x_ky_k|$, and consider the
  distance spheres $\partial B(x_k,|xz|)$ and $\partial B(y_k,|yz|)$; since all
  triangle inequalities among $x,y,z$ hold, these spheres are nonempty and
  intersect, and any $z_k$ in the intersection completes the comparison triangle.
  If any of the distances equals $\pi/\sqrt k$, Cheng's diameter theorem
  (Thm.\ 7.2.5 of the source) forces $M=S_k^n$, in which case there is nothing
  further to prove.
\end{proof}

\begin{proposition}[Prop. 12.2.3 — law of cosines in $S_k^n$]
  \label{prop:pet-ch12-law-of-cosines}
  \lean{PetersenLib.lawOfCosinesSpaceForm}
  Let a triangle in $S_k^n$ have side lengths $a,b,c$ and let $\alpha$ be the
  angle opposite $a$. Then
  \[
    k=0:\ a^2=b^2+c^2-2bc\cos\alpha;\qquad
    k=-1:\ \cosh a=\cosh b\cosh c-\sinh b\sinh c\cos\alpha;
  \]
  \[
    k=1:\ \cos a=\cos b\cos c+\sin b\sin c\cos\alpha.
  \]
  \source{petersen:page-0464,petersen:page-0465,petersen:page-0466}
\end{proposition}
\begin{proof}
  \uses{prop:pet-ch12-law-of-cosines}
  Let $p\in S_k^n$ and let $\sigma:[0,c]\to S_k^n$ be a unit-speed segment with
  $\sigma(0)$ at distance $b$ from $p$. Let $f_k$ be the modified distance
  function to $p$ (Cor.\ 4.3.4, Thm.\ 5.7.5 of the source), whose Hessian is
  $\mathrm{Hess}\,f_k=(1-kf_k)g$ (Ex.\ 4.3.2). Put $\rho=f_k\circ\sigma$, so
  $\dot\rho=g(\dot\sigma,\nabla f_k)$ and $\ddot\rho=1-k\rho$. Let $\alpha$ be the
  interior angle at $\sigma(0)$ between $\sigma$ and the segment joining $p$ to
  $\sigma(0)$, so $-\cos\alpha=g(\dot\sigma(0),\nabla r)$ where $r=|{\cdot}p|$.

  \emph{Case $k=0$.} Here $\rho=\tfrac12r^2\circ\sigma$, $\ddot\rho=1$, so
  $\rho(t)=\tfrac12b^2-b t\cos\alpha+\tfrac12t^2$. Setting $t=c$, $a=|p\sigma(c)|$
  gives $\tfrac12a^2=\tfrac12b^2-bc\cos\alpha+\tfrac12c^2$.

  \emph{Case $k=-1$.} Here $\rho=\cosh(r\circ\sigma)-1$, so
  $\ddot\rho-\rho=1$, with $\rho(0)=\cosh b-1$, $\dot\rho(0)=-\sinh b\cos\alpha$.
  The general solution is $\rho(t)=(\rho(0)+1)\cosh t+\dot\rho(0)\sinh t-1$;
  setting $t=c$, $a=|p\sigma(c)|$ gives
  $\cosh a-1=\cosh b\cosh c-\sinh b\sinh c\cos\alpha-1$.

  \emph{Case $k=1$.} Here $\rho=1-\cos(r\circ\sigma)$, so $\ddot\rho+\rho=1$,
  with $\rho(0)=1-\cos b$, $\dot\rho(0)=-\sin b\cos\alpha$. The general solution
  is $\rho(t)=(\rho(0)-1)\cos t+\dot\rho(0)\sin t+1$; setting $t=c$,
  $a=|p\sigma(c)|$ gives $1-\cos a=-\cos b\cos c-\sin b\sin c\cos\alpha+1$, i.e.\
  $\cos a=\cos b\cos c+\sin b\sin c\cos\alpha$.
\end{proof}

\begin{lemma}[Lem. 12.2.4 — Hessian comparison]
  \label{lem:pet-ch12-hessian-comparison}
  \lean{PetersenLib.hessianComparisonSupportSense}
  Let $(M,g)$ be complete, $p\in M$, $r(x)=|xp|$. If $\sec M\ge k$, then in the
  support sense everywhere $\mathrm{Hess}_kr\le(1-kt)g$, where $\mathrm{Hess}_k$
  denotes the Hessian of the modified distance function $f_k=f_k\circ r$.
  \source{petersen:page-0466,petersen:page-0467}
\end{lemma}
\begin{proof}
  \uses{lem:pet-ch12-hessian-comparison}
  This is Thm.\ 6.4.3 of the source when $r$ is smooth; the extension to points
  where $r$ is not smooth follows the same support-function argument as
  Lem.\ 7.1.9 of the source.
\end{proof}

\begin{theorem}[Thm. 12.2.2 — Toponogov comparison theorem, 1959]
  \label{thm:pet-ch12-toponogov}
  \lean{PetersenLib.toponogovComparisonTheorem}
  \uses{def:pet-ch12-hinge,def:pet-ch12-triangle,lem:pet-ch12-comparison-hinge-triangle-exist}
  Let $(M,g)$ be a complete Riemannian manifold with $\sec\ge k$.
  \begin{itemize}
    \item \emph{Hinge version.} For any hinge with vertices $p,x,y\in M$ forming
      angle $\alpha$ at $x$, and any comparison hinge in $S_k^n$ with vertices
      $p_k,x_k,y_k$ (\cref{lem:pet-ch12-comparison-hinge-triangle-exist}),
      $|py|\le|p_ky_k|$.
    \item \emph{Triangle version.} For any triangle in $M$, the interior angles
      are no smaller than the corresponding interior angles of a comparison
      triangle in $S_k^n$.
  \end{itemize}
  \source{petersen:page-0463,petersen:page-0464,petersen:page-0467,petersen:page-0468,petersen:page-0469,petersen:page-0470}
\end{theorem}
\begin{proof}
  \uses{thm:pet-ch12-toponogov,prop:pet-ch12-law-of-cosines,lem:pet-ch12-hessian-comparison}
  \emph{The hinge version implies the triangle version.} By
  \cref{prop:pet-ch12-law-of-cosines}, in $S_k^n$ increasing $|py|$ while
  holding $|px|,|xy|$ fixed strictly increases the angle at $x$; the hinge
  estimate applied to each pair of sides of a triangle in $M$, compared against
  the same side lengths in $S_k^n$, then yields angle comparisons on all three
  vertices.

  \emph{Hinge version.} Fix $p\in M$ and a geodesic $c:[0,L]\to M$, and select
  $\bar p\in S_k^n$, a segment $\bar c:[0,L]\to S_k^n$, with matching initial
  data: if $r=|{\cdot}p|$ is smooth at $c(0)$, take $|pc(0)|\le|\bar p\bar c(0)|$
  and $g(\nabla r,\dot c(0))\le g_k(\nabla \bar r,\dot{\bar c}(0))$; if $r$ is not
  smooth at $c(0)$, slide $c$ down a segment joining $p$ to $c(0)$ and pass to
  the limit, which lets one assume the strict inequality
  $|pc(0)|<|\bar p\bar c(0)|$ throughout. The claim is $|pc(t)|\le|\bar p\bar
  c(t)|$ for all $t\in[0,L]$.

  \emph{Case $k=0$.} Put $\rho(t)=\tfrac12(r\circ c(t))^2$,
  $\bar\rho(t)=\tfrac12(\bar r\circ\bar c(t))^2$. For small $t$ these are smooth
  with $\rho(0)<\bar\rho(0)$, $\dot\rho(0)\le\dot{\bar\rho}(0)$, and by
  \cref{lem:pet-ch12-hessian-comparison} (restricted to $\dot c,\dot{\bar c}$)
  $\ddot\rho=1\ge\ddot{\bar\rho}$ in the support sense. Hence
  $\psi=\bar\rho-\rho$ satisfies $\psi(0)>0$, $\dot\psi(0)\ge0$,
  $\ddot\psi\ge0$ in the support sense, so $\psi$ is convex, positive and
  increasing for small $t$, hence positive and increasing for all $t$; this
  gives $\bar\rho\ge\rho$, i.e.\ the hinge estimate.

  \emph{Case $k=-1$.} Put $\rho=\cosh(r\circ c)-1$,
  $\bar\rho=\cosh(\bar r\circ\bar c)-1$. As before $\rho(0)<\bar\rho(0)$,
  $\dot\rho(0)\le\dot{\bar\rho}(0)$, $\ddot\rho=\bar\rho+1$,
  $\ddot{\bar\rho}\le\rho+1$ in the support sense, so
  $\psi=\bar\rho-\rho$ satisfies $\psi(0)>0$, $\dot\psi(0)\ge0$,
  $\ddot\psi\ge\psi$ in the support sense. As long as $\psi>0$ it is convex; a
  convex function that starts positive and increasing cannot acquire a positive
  maximum, so $\psi$ keeps increasing and stays positive for all $t$.

  \emph{Case $k=1$.} Put $\rho=1-\cos(r\circ c)$,
  $\bar\rho=1-\cos(\bar r\circ\bar c)$, so $\psi=\bar\rho-\rho$ satisfies
  $\psi(0)>0$, $\dot\psi(0)\ge0$, $\ddot\psi\ge-\psi$ in the support sense. First
  perturb $c$ (assuming the interior angle at $c(0)$ positive — otherwise the
  hinge estimate is trivial) to arrange the strict inequality $\psi(0)>0$
  strictly. Compare $\psi$ to $\zeta$ solving
  $\ddot\zeta=-(1+\varepsilon)\zeta$, $\zeta(0)=\psi(0)>0$,
  $\dot\zeta(0)=\dot\psi(0)$. For small $t$,
  \[
    \frac{d^2}{dt^2}(\psi-\zeta)\ge-\psi+(1+\varepsilon)\zeta=(\zeta-\psi)+\varepsilon\zeta>0,
  \]
  so $\psi\ge\zeta$ for small $t$. On the interval
  $t<\big(\pi-\arctan(\dot\psi(0)\sqrt{1+\varepsilon}/\psi(0))\big)/\sqrt{1+\varepsilon}$
  where $\zeta>0$, suppose the ratio $\psi/\zeta$, which starts at $1$ and
  exceeds $1$ for small $t$, dips below $1$; it would then have a positive local
  maximum at some $t_0$. Using support functions $\psi_\delta\le\psi$ from below
  with $\ddot\psi_\delta\ge-\psi-\delta$, at such a maximum
  \[
    0\ge\frac{d^2}{dt^2}\Big(\frac{\psi_\delta}{\zeta}\Big)(t_0)
    \ge\frac{-\psi_\delta(t_0)-\delta}{\zeta(t_0)}+\frac{\psi_\delta(t_0)}{\zeta(t_0)}(1+\varepsilon)
    =\frac{\varepsilon\psi_\delta(t_0)-\delta}{\zeta(t_0)},
  \]
  using that the middle (first-derivative) term vanishes at a maximum of
  $\psi_\delta/\zeta$ and $\zeta$ solves the comparison ODE. As $\delta\to0$
  this becomes positive since $\psi_\delta(t_0)>0$, a contradiction. Hence
  $\psi\ge\zeta$ throughout the interval where $\zeta>0$. Letting
  $\varepsilon\to0$ and then $\psi(0)\to0$ (continuity) extends the estimate
  $\psi\ge0$ to all $t\le\pi$, which is the hinge estimate for $k=1$.
\end{proof}

\begin{remark}[Rem. 12.2.5 — role of segments in the proof]
  \label{rem:pet-ch12-toponogov-segment-remark}
  \lean{PetersenLib.toponogovSegmentRemark}
  \uses{thm:pet-ch12-toponogov}
  The proof never uses that the objects being compared in $M$ are segments; it
  only requires that the geodesics used in $S_k^n$ are segments. For $k\le0$
  this is automatic. For $k>0$ it means the comparison geodesic must have length
  $<\pi/\sqrt k$, the condition used in the $k=1$ case of the proof.
  \source{petersen:page-0470}
\end{remark}

\section{Sphere Theorems}

Throughout this section, after rescaling, one works with a closed Riemannian
$n$-manifold $(M,g)$ with $\sec\ge1$. Established consequences of the preceding
results (proved elsewhere in the source) include: $\mathrm{diam}(M,g)\le\pi$,
with equality only for $M=S^n(1)$; if $n$ is odd, $M$ is orientable; if $n$ is
even and $M$ orientable, $M$ is simply connected with
$\mathrm{inj}(M)\ge\pi/\sqrt{\max\sec}$; if $M$ is simply connected and
$\max\sec<4$, then $\mathrm{inj}(M)\ge\pi/\sqrt{\max\sec}$ and $M$ is homotopy
equivalent to a sphere.

\begin{theorem}[Thm. 12.3.1 — quarter-pinched sphere theorem, Rauch--Berger--Klingenberg]
  \label{thm:pet-ch12-quarter-pinched-sphere}
  \lean{PetersenLib.quarterPinchedSphereTheorem}
  \uses{thm:pet-ch12-toponogov}
  If $M$ is a simply connected closed Riemannian manifold with
  $1\le\sec\le4-\delta$, then $M$ is homeomorphic to a sphere.
  \source{petersen:page-0471}
\end{theorem}
\begin{proof}
  \uses{thm:pet-ch12-quarter-pinched-sphere,thm:pet-ch12-toponogov,def:pet-ch12-alpha-regular}
  The pinching gives $\mathrm{inj}(M)\ge\pi/\sqrt{4-\delta}$, so large discs
  exist around every point. Select $p,q\in M$ with $|pq|=\mathrm{diam}\,M$; then
  $\mathrm{diam}\,M\ge\mathrm{inj}\,M>\pi/2$. We show every $x\in M$ lies in
  $B(p,\pi/\sqrt{4-\delta})$ or $B(q,\pi/\sqrt{4-\delta})$, so $M$ is covered by
  two discs and hence homeomorphic to a sphere.

  Fix $x\in M$ and consider the triangle $p,x,q$. Suppose $|xq|>\pi/2$; we show
  $|px|<\pi/2$. Since $q$ maximizes distance from $p$, $q$ is not a regular
  point for the distance function to $p$, so for a segment $\overline{xq}$
  there is a segment $\overline{pq}$ with interior angle
  $\alpha=\angle(\overrightarrow{qx},\overrightarrow{qp})\le\pi/2$ at $q$. The
  hinge version of \cref{thm:pet-ch12-toponogov} gives
  \[
    \cos|px|\ge\cos|xq|\cos|pq|+\sin|xq|\sin|pq|\cos\alpha\ge\cos|xq|\cos|pq|.
  \]
  Since $|xq|,|pq|>\pi/2$, the right-hand side is positive, forcing
  $\cos|px|>0$, i.e.\ $|px|<\pi/2$. By symmetry the same argument with $p,q$
  interchanged shows that if $|px|>\pi/2$ then $|xq|<\pi/2$, so every $x$ lies
  in one of the two balls.
\end{proof}

\begin{theorem}[Thm. 12.3.2 — Berger 1962 and Grove--Shiohama diameter sphere theorem, 1977]
  \label{thm:pet-ch12-diameter-sphere}
  \lean{PetersenLib.diameterSphereTheorem}
  \uses{thm:pet-ch12-toponogov}
  If $(M,g)$ is closed with $\sec\ge1$ and $\mathrm{diam}>\pi/2$, then $M$ is
  homeomorphic to a sphere.
  \source{petersen:page-0472,petersen:page-0473}
\end{theorem}
\begin{proof}
  \uses{thm:pet-ch12-diameter-sphere,thm:pet-ch12-toponogov,def:pet-ch12-alpha-regular}
  \emph{Berger's index proof.} Select $p,q$ with $|pq|=\mathrm{diam}\,M>\pi/2$.
  It suffices to find a point at which all geodesic loops have length $>\pi$; we
  claim $p$ works. If not, there is a geodesic loop $c$ at $p$ of length
  $L(c)\le\pi$. Since $p$ maximizes distance from $q$, there is a segment
  $\overline{qp}$ forming a hinge with $c$ of interior angle $\le\pi/2$ at $p$,
  so by the hinge version of \cref{thm:pet-ch12-toponogov}
  \[
    0>\cos|pq|\ge\cos|pq|\cos L(c)+\sin|pq|\sin L(c)\cos\alpha\ge\cos|pq|\cos L(c),
  \]
  impossible unless $L(c)=0$. Having such a point, the conclusion follows as in
  the proof of Thm.\ 6.5.4 of the source.

  \emph{Grove--Shiohama proof.} Fix $p,q$ with $|pq|=\mathrm{diam}\,M>\pi/2$; the
  claim is that $|xp|$ has $q$ as its only critical point. Let
  $x\in M-\{p,q\}$ and let $\alpha$ be the interior angle at $x$ between
  segments $\overline{xp}$, $\overline{xq}$; suppose $\alpha\le\pi/2$. The hinge
  version gives
  \[
    0>\cos|pq|\ge\cos|px|\cos|xq|+\sin|px|\sin|xq|\cos\alpha\ge\cos|px|\cos|xq|,
  \]
  so $\cos|px|,\cos|xq|$ have opposite signs; if $\cos|px|>0$ then
  $\cos|pq|>\cos|xq|$, so $|xq|>|pq|=\mathrm{diam}\,M$, a contradiction (and
  symmetrically if $\cos|xq|>0$). Hence every $x\ne p,q$ has interior angle
  $>\pi/2$ between segments to $p$ and to $q$, so $x$ is a regular point for
  $|xp|$ unless $x=q$.

  Construct a vector field $X$ equal to the gradient of $|xp|$ near $p$ and the
  negative gradient of $|xq|$ near $q$, along which $|xp|$ increases everywhere;
  each $x\in M-\{p,q\}$ lies on a unique integral curve $c_x(t)$ of $X$. For $x$
  ranging over a small sphere $\partial B(p,\varepsilon)\cong S^{n-1}$, the
  integral curve reaches $\partial B(q,\varepsilon)\cong S^{n-1}$ at a time
  $t_x$ depending continuously (in fact smoothly) on $x$, giving a
  diffeomorphism
  \[
    \partial B(p,\varepsilon)\times[0,1]\to M-(B(p,\varepsilon)\cup B(q,\varepsilon)),
    \qquad(x,t)\mapsto c_x(t\cdot t_x).
  \]
  Gluing this to the two discs $B(p,\varepsilon)$, $B(q,\varepsilon)$ yields a
  continuous bijection $M\to S^n$, exhibiting $M$ as homeomorphic to a sphere
  (smoothness across the two boundary spheres is not guaranteed by this
  construction).
\end{proof}

\begin{theorem}[Thm. 12.3.3 — Brendle--Schoen 2008 and Petersen--Tao 2009]
  \label{thm:pet-ch12-cross-pinching}
  \lean{PetersenLib.crossPinchingClassification}
  Let $(M,g)$ be simply connected of dimension $n$. There is $\varepsilon(n)>0$
  such that if $1\le\sec\le4+\varepsilon$, then $M$ is diffeomorphic to a sphere
  or to one of the compact rank-one symmetric spaces (CROSS) $\mathbb{CP}^{n/2}$,
  $\mathbb{HP}^{n/4}$, $\mathbb{OP}^2$.
  \source{petersen:page-0474}
\end{theorem}

\begin{theorem}[Thm. 12.3.4 — Grove--Gromoll 1987 and Wilking 2001]
  \label{thm:pet-ch12-diameter-classification}
  \lean{PetersenLib.diameterHalfPiClassification}
  If $(M,g)$ is closed with $\sec\ge1$ and $\mathrm{diam}\ge\pi/2$, then one of
  the following holds:
  \begin{enumerate}
    \item $M$ is homeomorphic to a sphere;
    \item $M$ is isometric to a finite quotient $S^n(1)/\Gamma$ with $\Gamma$
      acting reducibly (an invariant subspace exists);
    \item $M$ is isometric to $\mathbb{CP}^m$, $\mathbb{HP}^m$, or
      $\mathbb{CP}^m/\mathbb{Z}_2$ for $m$ odd;
    \item $M$ is isometric to $\mathbb{OP}^2$.
  \end{enumerate}
  \source{petersen:page-0474}
\end{theorem}

\section{The Soul Theorem}

\begin{lemma}[Lem. 12.4.2 — Gromov's critical point estimate, 1981]
  \label{lem:pet-ch12-gromov-critical-point-estimate}
  \lean{PetersenLib.gromovCriticalPointEstimate}
  \uses{thm:pet-ch12-toponogov,def:pet-ch12-alpha-regular}
  If $(M,g)$ is a complete open manifold with $\sec\ge0$, then for every $p\in M$
  the distance function $|xp|$ has no critical points outside some ball
  $B(p,R)$. In particular $M$ has the topology of a compact manifold with
  boundary.
  \source{petersen:page-0475}
\end{lemma}
\begin{proof}
  \uses{lem:pet-ch12-gromov-critical-point-estimate,thm:pet-ch12-toponogov}
  Let $x$ be critical for $|xp|$ and let $y$ satisfy
  $\angle(\overrightarrow{px},\overrightarrow{py})\le\pi/3$. The hinge version
  of \cref{thm:pet-ch12-toponogov} (case $k=0$) applied at $p$ gives
  \[
    |xy|^2\le|py|^2+|xp|^2-2|py||xp|\cos\tfrac\pi3=|py|^2+|xp|^2-|py||xp|.
  \]
  Since $x$ is critical for $p$, choose segments $\overline{px}$,
  $\overline{xy}$ with angle $\le\pi/2$ at $x$; the hinge version at $x$ gives
  \[
    |py|^2\le|xp|^2+|xy|^2\le2|xp|^2+|py|^2-|py||xp|,
  \]
  which forces $|py|\le2|xp|$. There are at most a fixed number
  $\le v(n-1,1,\pi)/v(n-1,0,\pi/6)\le6^{n-1}$ of unit vectors at $p$ pairwise
  making angle $>\pi/3$ (disjoint balls of radius $\pi/6$ around them embed in
  the unit sphere). Consequently there are at most $6^{n-1}$ critical points
  $x_i$ for $|xp|$ with $|x_{i+1}p|>2|x_ip|$ along any such chain, so no critical
  points occur outside some ball $B(p,R)$.
\end{proof}

\begin{definition}[totally convex]
  \label{def:pet-ch12-totally-convex}
  \lean{PetersenLib.IsTotallyConvex}
  A subset $A\subset M$ of a Riemannian manifold is \emph{totally convex} if
  every geodesic in $M$ joining two points of $A$ lies entirely in $A$.
  \source{petersen:page-0476}
\end{definition}

\begin{example}[Ex. 12.4.3 — souls of the flat cylinder]
  \label{ex:pet-ch12-totally-convex-cylinder}
  \lean{PetersenLib.flatCylinderSouls}
  \uses{def:pet-ch12-totally-convex}
  On the flat cylinder $\mathbb{R}\times S^1$, every circle $\{p\}\times S^1$ is
  a geodesic and totally convex; no point of $M$ is totally convex, and every
  such circle is a soul.
  \source{petersen:page-0476}
\end{example}

\begin{example}[Ex. 12.4.4 — soul of a paraboloid of revolution]
  \label{ex:pet-ch12-totally-convex-paraboloid}
  \lean{PetersenLib.paraboloidSoul}
  \uses{def:pet-ch12-totally-convex}
  Let $(M,g)$ be a smooth rotationally symmetric metric
  $dr^2+\rho^2(r)d\theta^2$ on $\mathbb{R}^2$ with $\rho'<0$ (a paraboloid of
  revolution). Every geodesic from the origin $r=0$ is a ray to infinity, so the
  origin is a totally convex soul; most other points admit geodesic loops.
  \source{petersen:page-0476}
\end{example}

\begin{lemma}[Lem. 12.4.5 — superlevel sets of concave functions]
  \label{lem:pet-ch12-concave-superlevel-convex}
  \lean{PetersenLib.concaveSuperlevelTotallyConvex}
  \uses{def:pet-ch12-totally-convex}
  If $f:(M,g)\to\mathbb{R}$ is concave, meaning $\mathrm{Hess}\,f\le0$ weakly
  everywhere, then every superlevel set $A=\{x\in M\mid f(x)\ge a\}$ is totally
  convex.
  \source{petersen:page-0476}
\end{lemma}
\begin{proof}
  \uses{lem:pet-ch12-concave-superlevel-convex}
  For a geodesic $c$ in $M$, $f\circ c$ has nonpositive weak second derivative,
  hence is concave on $\mathbb{R}$; the minimum of a concave function on any
  compact interval is attained at an endpoint. If both endpoints of a segment
  in $M$ lie in $A$, i.e.\ $f\ge a$ there, this shows $f\ge a$ along the whole
  segment, so the segment lies in $A$.
\end{proof}

\begin{lemma}[Lem. 12.4.6 — properness and concavity of the Busemann infimum]
  \label{lem:pet-ch12-busemann-proper-concave}
  \lean{PetersenLib.busemannInfimumProperConcave}
  \uses{def:pet-ch12-totally-convex,lem:pet-ch12-concave-superlevel-convex}
  Let $(M,g)$ be complete, noncompact, $\sec\ge0$, and $p\in M$. Let
  $R_p=\{c:[0,\infty)\to M\mid c(0)=p\}$ be the set of rays from $p$, and set
  $f=\inf_{c\in R_p}b_c$, where $b_c$ is the Busemann function of $c$. Then $f$
  is proper and concave.
  \source{petersen:page-0477}
\end{lemma}
\begin{proof}
  \uses{lem:pet-ch12-busemann-proper-concave,lem:pet-ch12-concave-superlevel-convex}
  \emph{Concavity of Busemann functions.} In nonnegative Ricci curvature
  Busemann functions are superharmonic (§7.3.2 of the source); the same
  Hessian-estimate argument, using that $r(x)=|xy|$ satisfies
  $\mathrm{Hess}\,r=0$ radially and $\mathrm{Hess}\,r\le r^{-1}g$ transverse to
  the radial direction at smooth points (and hence in the support sense
  everywhere, by the usual extension to nonsmooth points), shows that each
  $b_c$ has nonpositive Hessian in the weak sense, i.e.\ is concave.

  \emph{Concavity of $f$.} An infimum of concave functions is concave, so $f$ is
  concave.

  \emph{Properness.} Suppose some superlevel set
  $A=\{x\in M\mid f(x)\ge a\}$ is noncompact; since larger superlevel sets are
  contained in smaller ones we may take $a\le0$. As $f(p)=0$ (every $b_c$
  vanishes at $p$), $p\in A$. Take $p_i\in A$ with $p_i\to\infty$; passing to a
  subsequence, the segments $\overline{pp_i}$ converge to a ray $c\in R_p$ (the
  standard construction of a ray as a limit of segments). By
  \cref{lem:pet-ch12-concave-superlevel-convex}, $A$ is totally convex, so each
  segment $\overline{pp_i}$, and hence in the limit the ray $c$, lies in the
  closed set $A$. But then
  \[
    a\le f(c(t))\le b_c(c(t))=-t\to-\infty,
  \]
  a contradiction. Hence every superlevel set of $f$ is compact, i.e.\ $f$ is
  proper.
\end{proof}

\begin{lemma}[Lem. 12.4.7 — interior, boundary, and the supporting hyperplane property]
  \label{lem:pet-ch12-totally-convex-boundary-structure}
  \lean{PetersenLib.totallyConvexBoundaryStructure}
  \uses{def:pet-ch12-totally-convex}
  If $A\subset(M,g)$ is closed and totally convex, then $A$ has an interior
  $\mathrm{int}\,A$ and a boundary $\partial A=A-\mathrm{int}\,A$, where
  $\mathrm{int}\,A$ is a totally convex submanifold of $M$ with closure $A$.
  Moreover for each $x\in\partial A$ there is an inward vector
  $w\in T_xM$ such that every segment $\overline{xy}$ with $y\in\mathrm{int}\,A$
  satisfies $\angle(w,\overline{xy})<\pi/2$ (the \emph{supporting hyperplane
  property}); in particular the distance function to a subset of
  $\mathrm{int}\,A$ has no critical points on $\partial A$.
  \source{petersen:page-0477,petersen:page-0478,petersen:page-0479,petersen:page-0480}
\end{lemma}
\begin{proof}
  \uses{lem:pet-ch12-totally-convex-boundary-structure}
  Let $\varepsilon(p):M\to(0,\infty)$ be the convexity radius function of
  Thm.\ 6.4.8 of the source, so that $r_p=|{\cdot}p|$ is smooth and strictly
  convex on $B(p,\varepsilon(p))-\{p\}$.

  \emph{Identifying $N=\mathrm{int}\,A$.} Let $k$ be the largest integer such
  that $A$ contains a $k$-dimensional submanifold of $M$ (so $k=0$ forces $A$ to
  be a single point, since two points of $A$ would force a segment, hence a
  $1$-dimensional submanifold, into $A$). Let $N\subset A$ be the union of all
  $k$-dimensional submanifolds of $M$ contained in $A$. To see $N$ is an
  embedded $k$-manifold, fix $p\in N$ with $N_p\subset A$ a $k$-dimensional
  embedded submanifold through $p$, and $\delta\in(0,\varepsilon(p))$,
  $\delta<\mathrm{inj}_p$, with $B(p,\delta)\cap N_p=N_p$. If some
  $q\in A\cap B(p,\delta)-N_p$ existed, joining every point of
  $B(p,\delta)\cap N_p$ to $q$ by its unique segment produces, away from $q$, a
  $(k+1)$-dimensional submanifold in $A$, contradicting maximality of $k$; hence
  $B(p,\delta)\cap A=N_p$, showing $N$ is embedded near $p$. The same argument
  (using a $(k-1)$-dimensional submanifold of $N$ through $p$ perpendicular to a
  segment to $q\in A$ with $|pq|<\mathrm{inj}_p$, and coning over it towards
  $q$) shows every such segment $\overline{pq}$ lies in $N$ except possibly at
  $q$; hence $N$ is dense in $A$, justifying $\mathrm{int}\,A:=N$,
  $\partial A:=A-N$.

  \emph{Supporting hyperplane property.} Since $N$ is totally geodesic, its
  tangent spaces are preserved by parallel translation along curves in $N$, so
  for $p\in\partial A$ one obtains a well-defined $k$-plane $T_pA\subset T_pM$
  by parallel transport from $N$ along curves ending at $p$. Define the tangent
  cone
  \[
    C_pA=\{v\in T_pM\mid\exp_p(tv)\in N\text{ for some (hence all small) }t>0\}
    \subset T_pA,
  \]
  which is open in $T_pA$. If $C_pA$ contained antipodal vectors $\pm v$ then
  short segments through $p$ with both endpoints in $N$ would force $p\in N$,
  contradiction; so $C_pA$ contains no antipodal pair.

  For $p\in\partial A$ in the (dense) set of boundary points admitting
  $q\in A$ with $|pq|=\varepsilon$ realizing $\varepsilon=\min\{|px|\mid x\in
  A,\,x\ne p\}$ for small $\varepsilon$ so that $r_q=|{\cdot}q|$ is smooth and
  convex near $p$: set $A'=A\cap\overline B(q,\varepsilon)$,
  $N'=A\cap B(q,\varepsilon)\subset N$, so $p\in\partial A'$,
  $C_pA'\subset C_pA$, $T_pA=T_pA'$. Since $r_q$ is smooth at $p$,
  $C_p\overline B(q,\varepsilon)=\{v\in T_pM\mid\angle(v,-\nabla
  r_q)<\pi/2\}$, so
  $C_pA'=\{v\in T_pA\mid\angle(v,-\nabla r_q)<\pi/2\}$. If this were a strict
  subset of $C_pA$, openness of $C_pA$ in $T_pA$ would force $C_pA$ to contain
  antipodal vectors, contradiction; so $C_pA=C_pA'$ is exactly the open half
  space $\{v\in T_pA\mid\angle(v,-\nabla r_q)<\pi/2\}$, giving the supporting
  hyperplane property with $w=-\nabla r_q$ at such $p$.

  For general $p\in\partial A$, approximate by $p_i\to p$ of the above type,
  with $C_{p_i}A=\{v\in T_{p_i}A\mid\angle(v,-w_i)<\pi/2\}$; these half-spaces
  accumulate on a half space $\{v\in T_pA\mid\angle(v,-w)<\pi/2\}$ for some
  limit $w$, and by continuity $C_pA\subset\{v\in T_pA\mid\angle(v,-w)\le\pi/2\}$;
  since $C_pA$ is open in $T_pA$ it is contained in the open half space, giving
  the supporting hyperplane property at $p$ with this $w$. The stated
  non-criticality of the distance function to a subset of $\mathrm{int}\,A$ on
  $\partial A$ is immediate from this property, since any inward direction
  $\overline{xy}$, $y\in\mathrm{int}\,A$, then makes angle $<\pi/2$ with the
  supporting direction $w$, so the distance function cannot be minimized (nor
  have vanishing gradient) at $x$.
\end{proof}

\begin{lemma}[Lem. 12.4.8 — concavity and uniqueness of distance to the boundary]
  \label{lem:pet-ch12-distance-to-boundary-concave}
  \lean{PetersenLib.distanceToBoundaryConcave}
  \uses{def:pet-ch12-totally-convex,lem:pet-ch12-totally-convex-boundary-structure}
  Let $(M,g)$ have $\sec\ge0$ and let $A\subset M$ be totally convex. The
  function $r:A\to\mathbb{R}$, $r(x)=|x\partial A|$, is concave on $A$. When
  $\sec>0$, any maximum of $r$ is unique.
  \source{petersen:page-0480,petersen:page-0481}
\end{lemma}
\begin{proof}
  \uses{lem:pet-ch12-distance-to-boundary-concave,lem:pet-ch12-totally-convex-boundary-structure}
  Fix $q\in\mathrm{int}\,A$ and $p\in\partial A$ with $|pq|=r(q)$, and a segment
  $\overline{pq}\subset A$. In exponential coordinates at $p$ let $H$ be the
  image of the hyperplane perpendicular to $\overline{pq}$, so $H$ is
  perpendicular to $\overline{pq}$, has vanishing second fundamental form at
  $p$, and $H\cap\mathrm{int}\,A=\varnothing$ (by the supporting hyperplane
  property of \cref{lem:pet-ch12-totally-convex-boundary-structure}). Then
  $f(x)=|xH|$ is a support function from above for $r$ along $\overline{pq}$.

  Fix $p'\ne p,q$ on $\overline{pq}$, so $f$ is smooth at $p'$ (as in §5.7.3 of
  the source, unless $p'=q$). Since $\overline{pq}$ is an integral curve of
  $\nabla f$, evaluating the fundamental (Riccati) equation of §3.2.5 of the
  source on a parallel field $E$ along $\overline{pq}$ tangent to $H$ at $p$
  gives
  \[
    \frac{d}{dt}\mathrm{Hess}\,f(E,E)=-R(E,\nabla f,\nabla f,E)-\mathrm{Hess}^2f(E,E)\le0,
  \]
  and since $\mathrm{Hess}\,f(E,E)=0$ at $p$, $\mathrm{Hess}\,f(E,E)\le0$ (and
  $<0$ if $\sec>0$) along $\overline{pq}$, so $f$ is a smooth support function
  for $r$ on an open dense subset of $\overline{pq}$.

  If $f$ is not smooth at $q$, take instead the hypersurface $H'$ perpendicular
  to $\overline{p'q}$ at $p'$ with vanishing second fundamental form there; for
  $p'$ close to $q$, $|{\cdot}H'|$ is smooth at $q$ with nonpositive (negative)
  Hessian there, and $|pp'|+|{\cdot}H'|$ is a support function for $r$ at $q$:
  the two functions agree at $q$, and for $x$ with $r(x)>|pp'|$ one finds
  $z\in H'$ with $r(x)=r(z)+|xH'|$, so it suffices that $r(z)\le|pp'|$ for
  $z\in H'$. Since $f$ (hence $r$) is concave near $p'$, and a segment
  $\overline{p'z}\subset H'$ is perpendicular to $\overline{p'q}$, concavity of
  $r$ along $\overline{p'z}$ gives $r(z)\le r(p')=|pp'|$.

  Finally, for $\phi:[0,\infty)\to[0,\infty)$ concave with $\phi(0)=0$,
  $\phi'>0$, $\phi\circ r$ is concave (strictly, if $\phi$ is strictly concave
  and $\sec>0$); a strictly concave function has at most one maximum, by the
  same uniqueness argument as for the center of mass (§6.2.2 of the source).
\end{proof}

\begin{theorem}[Thm. 12.4.1 — soul theorem, Gromoll--Meyer 1969 and Cheeger--Gromoll 1972]
  \label{thm:pet-ch12-soul-theorem}
  \lean{PetersenLib.soulTheorem}
  \uses{lem:pet-ch12-busemann-proper-concave,lem:pet-ch12-distance-to-boundary-concave,cor:pet-ch12-normal-bundle-diffeomorphism}
  If $(M,g)$ is a complete noncompact Riemannian manifold with $\sec\ge0$, then
  $M$ contains a soul $S\subset M$: a closed totally convex submanifold such
  that $M$ is diffeomorphic to the normal bundle of $S$. When $\sec>0$, the soul
  is a point and $M$ is diffeomorphic to $\mathbb{R}^n$.
  \source{petersen:page-0475,petersen:page-0481,petersen:page-0482}
\end{theorem}
\begin{proof}
  \uses{thm:pet-ch12-soul-theorem,lem:pet-ch12-busemann-proper-concave,lem:pet-ch12-concave-superlevel-convex,lem:pet-ch12-distance-to-boundary-concave,lem:pet-ch12-totally-convex-boundary-structure,cor:pet-ch12-normal-bundle-diffeomorphism}
  Let $f$ be the proper concave function of
  \cref{lem:pet-ch12-busemann-proper-concave}, and put
  $C_1=\{x\in M\mid f(x)=\max f\}$. Since $f$ is proper and concave, $C_1$ is
  nonempty, compact, and (as a superlevel set,
  \cref{lem:pet-ch12-concave-superlevel-convex}) totally convex. If $\sec>0$,
  \cref{lem:pet-ch12-distance-to-boundary-concave} applied to superlevel sets
  $A=\{f\ge a\}$, whose boundary is $f^{-1}(a)$ and on which $f(x)=|x\partial A|$
  up to an additive constant, shows the maximum of such a distance function is
  unique, so $C_1$ is a single point.

  If $C_1$ is a submanifold (with empty boundary as a totally convex set), stop
  and set $S=C_1$: every point of $M$ lies on the boundary of some superlevel
  set of $f$, so by the supporting hyperplane property of
  \cref{lem:pet-ch12-totally-convex-boundary-structure} the distance function
  $|xC_1|$ has no critical points on $M-C_1$.

  Otherwise $C_1$ is a compact totally convex set with nonempty boundary, and by
  \cref{lem:pet-ch12-distance-to-boundary-concave} the distance function
  $|x\partial C_1|$ is concave on $C_1$; let
  $C_2=\{x\in C_1\mid|x\partial C_1|=\max\}$, again nonempty, compact, totally
  convex, with a unique point if $\sec>0$. Iterating produces
  $C_1\supset C_2\supset\cdots$, terminating in at most $n=\dim M$ steps at a
  submanifold (with empty boundary as a totally convex subset) $S=C_k$: indeed
  if $\dim C_i=\dim C_{i+1}$, then $\mathrm{int}\,C_{i+1}$ is open in
  $\mathrm{int}\,C_i$, so for $p\in\mathrm{int}\,C_{i+1}$ there is $\delta$ with
  $B(p,\delta)\cap\mathrm{int}\,C_{i+1}=B(p,\delta)\cap\mathrm{int}\,C_i$; a
  segment $c$ from $p$ to $\partial C_i$ has $|x\partial C_i|$ strictly
  increasing along it near $p$, yet $c$ runs through
  $B(p,\delta)\cap\mathrm{int}\,C_i=B(p,\delta)\cap\mathrm{int}\,C_{i+1}$ on
  which $|x\partial C_i|$ is by construction constant (equal to its max on
  $C_i$) near $p$, a contradiction; hence $\dim C_i>\dim C_{i+1}$ strictly at
  each step until a step is a submanifold with empty boundary.

  Set $S=C_k$. At every stage, any point outside $C_{i+1}$ lies on the boundary
  of a superlevel set of $f$ (for $i=0$) or of $|x\partial C_i|$ ($i\ge1$), so
  by the supporting hyperplane property
  (\cref{lem:pet-ch12-totally-convex-boundary-structure}) the distance function
  $|xS|$ has no critical points on $M-S$. By
  \cref{cor:pet-ch12-normal-bundle-diffeomorphism}, $M$ is diffeomorphic to the
  normal bundle of $S$; if $\sec>0$, $S$ is a point (by the uniqueness of
  maxima established above at every stage) and the normal bundle of a point is
  $\mathbb{R}^n$.
\end{proof}

\begin{theorem}[Thm. 12.4.9 — Sharafutdinov 1978 and Perel'man 1993]
  \label{thm:pet-ch12-soul-uniqueness-submetry}
  \lean{PetersenLib.sharafutdinovSubmetryToSoul}
  \uses{thm:pet-ch12-soul-theorem}
  Let $S$ be a soul of a complete Riemannian manifold with $\sec\ge0$, arising
  from the construction of \cref{thm:pet-ch12-soul-theorem}.
  \begin{enumerate}
    \item (Sharafutdinov) There is a distance-nonincreasing map
      $\mathrm{Sh}:M\to S$ with $\mathrm{Sh}|_S=\mathrm{id}$; in particular all
      souls of $M$ are isometric to each other.
    \item (Perel'man) $\mathrm{Sh}$ is a submetry; consequently $S$ is a point
      whenever all sectional curvatures based at some single point of $M$ are
      positive.
  \end{enumerate}
  \source{petersen:page-0482}
\end{theorem}

\section{Finiteness of Betti Numbers}

\begin{lemma}[Lem. 12.5.2 — bounded corank]
  \label{lem:pet-ch12-corank-bound}
  \lean{PetersenLib.corankBound}
  \uses{thm:pet-ch12-toponogov,def:pet-ch12-alpha-regular}
  Fix a Riemannian $n$-manifold $M$ with $\sec\ge0$. The \emph{corank} of
  $A\subset M$ is the largest $k$ such that there are balls
  $B(p_1,r_1),\dots,B(p_k,r_k)$ with: (a) a critical point $x_i$ for $p_i$ with
  $|p_ix_i|=10r_i$; (b) $r_i\ge3r_{i-1}$ for $i\ge2$; (c)
  $A\subset\bigcap_iB(p_i,r_i)$. The corank of any $A\subset M$ is bounded by
  $100^n$.
  \source{petersen:page-0485,petersen:page-0486,petersen:page-0487}
\end{lemma}
\begin{proof}
  \uses{lem:pet-ch12-corank-bound,thm:pet-ch12-toponogov}
  Suppose $A$ has corank $k>100^n$, witnessed by
  $B(p_1,r_1),\dots,B(p_k,r_k)$ and critical points $x_1,\dots,x_k$. Fix
  $z\in A$. Since $v(n-1,1,\pi)/v(n-1,1,1/12)\le(12\pi)^{n-1}\le40^n<k$, two of
  the segments $\overline{zx_i}$, $\overline{zx_j}$ ($i<j$) form an angle
  $<1/6$ at $z$. From $|zp_i|\le r_i$, $|zp_j|\le r_j$ and $r_j\ge3r_i$,
  \[
    |zx_i|\le10r_i+|zp_i|\le11r_i,\qquad|zx_j|\ge10r_j-r_j\ge9r_j>|zx_i|.
  \]
  The hinge version of \cref{thm:pet-ch12-toponogov} at $z$ gives
  \[
    |x_ix_j|^2\le|zx_j|^2+|zx_i|^2-2|zx_i||zx_j|\cos\tfrac16
    =|zx_j|^2+|zx_i|^2-\tfrac{31}{16}|zx_i||zx_j|\le\big(|zx_j|-\tfrac34|zx_i|\big)^2,
  \]
  i.e.\ $|x_ix_j|\le|zx_j|-\tfrac34|zx_i|$. The triangle inequality gives
  \[
    |p_ix_j|\ge|zx_j|-|zp_i|\ge10r_j-|zp_j|-|zp_i|\ge8r_j\ge24r_i\ge2|p_ix_i|,
  \]
  hence $|x_ix_j|\ge|p_ix_j|$. Since $x_i$ is critical for $p_i$, the hinge
  version at $x_i$ gives
  \[
    |p_ix_j|^2\le|p_ix_i|^2+|x_ix_j|^2\le\big(|x_ix_j|+\tfrac12|p_ix_i|\big)^2,
  \]
  so $|p_ix_j|\le|x_ix_j|+5r_i$, and then
  $|zx_j|\le|p_ix_j|+|zp_i|\le|x_ix_j|+6r_i$. On the other hand
  $|zx_i|\ge10r_i-|zp_i|\ge9r_i$, so
  $|x_ix_j|\le|zx_j|-\tfrac34|zx_i|\le|zx_j|-\tfrac{27}4r_i$, giving
  \[
    |x_ix_j|+\tfrac{27}4r_i\le|zx_j|\le|x_ix_j|+6r_i,
  \]
  a contradiction since $27/4>6$. Hence the corank is at most $100^n$.
\end{proof}

\begin{lemma}[Lem. 12.5.3 — Mayer--Vietoris rank inequality]
  \label{lem:pet-ch12-mayer-vietoris-rank}
  \lean{PetersenLib.mayerVietorisRankInequality}
  Let $B_i^j$, $i=1,\dots,m$, $j=0,\dots,n+1$, be bounded open subsets of an
  $n$-manifold $M$ with $\overline{B_i^j}\subset B_i^{j+1}$. Then
  \[
    \mathrm{rank}_*\Big(\bigcup_iB_i^0\subset\bigcup_iB_i^{n+1}\Big)
    \le\sum_{l=0}^{n}\sum_{i_0<\cdots<i_l}
    \mathrm{rank}_*\Big(\bigcap_{s=0}^lB_{i_s}^0\subset\bigcap_{s=0}^lB_{i_s}^{n+1-l}\Big).
  \]
  \source{petersen:page-0488,petersen:page-0489,petersen:page-0490}
\end{lemma}
\begin{proof}
  \uses{lem:pet-ch12-mayer-vietoris-rank}
  Form the generalized Mayer--Vietoris double complex of singular chains with
  coefficients $\mathbb{F}$: $C_{p,q}^j=\bigoplus_{i_0<\cdots<i_q}C_p\big(\bigcap_{s=0}^qB_{i_s}^j\big)$,
  with the usual boundary $\partial:C_{p,q}^j\to C_{p-1,q}^j$ and the
  Mayer--Vietoris differential $\delta:C_{p,q}^j\to C_{p,q-1}^j$ induced (with
  alternating signs) by the inclusions
  $C_p\big(\bigcap_{s=0}^qB_{i_s}^j\big)\to C_p\big(\bigcap_{s\ne l}B_{i_s}^j\big)$;
  these satisfy $\partial\delta=\delta\partial$, so
  $\bar\partial:=(-1)^p\delta$ anticommutes with $\partial$, and
  $D=\partial+\bar\partial$ is a boundary map on
  $\bigoplus_{l=0}^kC^j_{l,k-l}$. Augment by the natural chain map
  $\bar\partial:C^j_{p,0}\to C_p\big(\bigcup_iB_i^j\big)$, and set
  $C_p^j=\mathrm{im}\,\bar\partial$. Any $D$-cycle $(c_{l,k-l}^j)$ gives a
  $\bar\partial$-cycle $\bar\partial c_{k,0}^j$, since
  $\partial\bar\partial c_{k,0}^j=-\bar\partial\partial c_{k,0}^j=\bar\partial\delta
  c_{k-1,1}^j=0$; conversely every $\bar\partial$-cycle lifts to a $D$-cycle by
  successively solving $\partial c_{l,k-l}^j=-\bar\partial c_{l-1,k-l+1}^j$
  using exactness of the $\bar\partial$-columns, so $D$-homology is isomorphic
  to $\bar\partial$-homology, i.e.\ to the ordinary singular homology of
  $\bigcup_iB_i^j$; this correspondence represents homology classes of
  $\bigcup_iB_i^j$ by chains supported in the pairwise intersections
  $\bigcap_sB_{i_s}^j$.

  Fix a finite-dimensional subspace $Z_k\subset\bigoplus_{l=0}^kC^0_{l,k-l}$ of
  $D$-cycles containing no nontrivial $D$-boundary, and let
  $Z_k^0=\bar\partial Z_k\subset C_k^0$ be the isomorphic image in
  $\bar\partial$-cycles. We construct a filtration
  $Z_k^0\supset Z_k^1\supset\cdots\supset Z_k^{k+1}$ together with linear maps
  $L^l:Z_k^l\to C^0_{l,k-l}$ whose images consist of $\bar\partial$-cycles,
  inductively: set $L^0(z)=z_{0,k}$ (trivially $\partial z_{0,k}=0$) and
  $Z_k^1=\{z\in Z_k^0\mid f^1(z_{0,k})\in\partial C^1_{1,k}\}$, where
  $f^l:C_*(\bigcup_iB_i^{l-1})\to C_*(\bigcup_iB_i^l)$ is induced by inclusion.
  Given $L^l$ with $\bar\partial L^l(z)=f^l\cdots f^1(\bar\partial z_{l,k-l})$
  and $\partial L^l(z)=0$, choose inverses
  $\bar\partial^{-1}:\partial C^j_{p,q}\to C^j_{p,q}$ with
  $\bar\partial\bar\partial^{-1}=\mathrm{id}$, set
  $Z_k^{l+1}=\{z\in Z_k^l\mid f^{l+1}L^l(z)\in\partial C^{l+1}_{k+1,k-l}\}$, and
  \[
    L^{l+1}(z)=f^{l+1}\cdots f^1(z_{l+1,k-l-1})+\bar\partial^{-1}f^{l+1}L^l(z),
  \]
  which one checks (using that $(z_{l,k-l})$ is a $D$-cycle) again satisfies the
  same two identities with $l$ replaced by $l+1$. At each step, restriction
  $z\mapsto z_{l,k-l}$ is linear and factors through the rank of
  $\bigcap_{i_0<\cdots<i_{k-l}}B_{i_s}^l\subset\bigcap_sB_{i_s}^{l+1}$, giving
  \[
    \dim(Z_k^l/Z_k^{l+1})\le\sum_{i_0<\cdots<i_{k-l}}
    \mathrm{rank}_*\Big(\bigcap_{s=0}^{k-l}B_{i_s}^l\subset\bigcap_{s=0}^{k-l}B_{i_s}^{l+1}\Big).
  \]
  Finally, when $z\in Z_k^{k+1}$, $f^{k+1}L^k(z)$ is a $\partial$-boundary and
  $\bar\partial f^{k+1}L^k(z)=f^{k+1}\cdots f^1(z)$, so
  $f^{k+1}\cdots f^1(z)\in C_k^{k+1}$ is a $\bar\partial$-boundary. Choosing
  $Z_k^0$ to map isomorphically onto its image in
  $H_k\big(\bigcup_iB_i^{k+1}\big)$ (possible since $Z_k^0$ was chosen with no
  nontrivial $D$-boundary), the rank of
  $\mathrm{rank}_*(\bigcup_iB_i^0\subset\bigcup_iB_i^{n+1})$ in degree $k$ is
  bounded by $\dim Z_k^0=\sum_{l=0}^k\dim(Z_k^l/Z_k^{l+1})+\dim Z_k^{k+1}$,
  and summing over $k=0,\dots,n$ and using $\dim Z_k^{k+1}=0$ (since such classes
  vanish in $\bigcup_iB_i^{n+1}$-homology) gives the stated bound.
\end{proof}

\begin{lemma}[Lem. 12.5.4 — content ratio bound]
  \label{lem:pet-ch12-content-ratio-bound}
  \lean{PetersenLib.contentRatioBound}
  \uses{lem:pet-ch12-corank-bound,lem:pet-ch12-mayer-vietoris-rank}
  Let $M$ be a Riemannian $n$-manifold with $\sec\ge0$, let
  $\mathrm{cont}\,B(p,r)=\mathrm{rank}_*(B(p,r/2)\subset B(p,r))$ be the
  \emph{content} of a ball, and let $\mathcal{C}(k)$ be the largest content of
  any ball of corank $\ge k$ (\cref{lem:pet-ch12-corank-bound}). There is
  $C(n)$, depending only on $n$, with $\mathcal{C}(k)\le
  C(n)\,\mathcal{C}(k+1)$.
  \source{petersen:page-0491,petersen:page-0492}
\end{lemma}
\begin{proof}
  \uses{lem:pet-ch12-content-ratio-bound,lem:pet-ch12-corank-bound,lem:pet-ch12-mayer-vietoris-rank}
  Call a ball $B(p,R)$ \emph{compressible} if it contains a ball $B(x,r)$ with
  $r\le R/2$ and $\mathrm{cont}\,B(x,r)\ge\mathrm{cont}\,B(p,R)$; otherwise
  \emph{incompressible}. $\mathcal{C}(k)$ is realized by some incompressible
  $B(p,R)\in\mathcal{B}(k)$. For $x\in B(p,R/4)$, $r\le R/100$, consider the
  chain $B(x,R/10)\subset B(p,R/5)\subset B(x,R/2)\subset B(p,R)$. If $x$ had no
  critical point in $B(x,R/2)-B(x,R/10)$, then
  $\mathrm{rank}_*(B(x,R/10)\subset B(x,R/2))=\mathrm{rank}_*(B(p,R/5)\subset
  B(x,R/2))\ge\mathrm{rank}_*(B(p,R/5)\subset B(p,R))$, giving
  $\mathrm{cont}\,B(p,R)\le\mathrm{cont}\,B(x,R/2)$, contradicting
  incompressibility. So there is a critical point $y$ for $x$ in
  $B(x,R/2)-B(x,R/10)$; adjoining $B(x,|xy|/10)$ to the balls witnessing
  $B(p,R)\in\mathcal{B}(k)$ shows $B(x,r)\in\mathcal{B}(k+1)$ for every such $x,r$.

  Cover $B(p,R/5)$ by balls $B(p_i,10^{-n-3}R)$, $i=1,\dots,m$, with the balls
  $B(p_i,10^{-n-3}R/2)$ pairwise disjoint, so
  $m\le v(n,0,R)/v(n,0,10^{-n-3}R/2)=2^n\cdot10^{n(n+3)}$. Since
  $B(p_i,R/2)\subset B(p,R)$,
  \[
    \mathrm{cont}\,B(p,R)\le\mathrm{rank}_*\Big(\bigcup_iB(p_i,10^{-n-3}R)\subset\bigcup_iB(p_i,R/2)\Big).
  \]
  Applying \cref{lem:pet-ch12-mayer-vietoris-rank} to the family
  $B_i^j=B(p_i,10^{j-n-3}R)$, $j=0,\dots,n+1$, and using that any nonempty
  intersection $\bigcap_{\ell=0}^sB(p_{i_\ell},10^{j-n-3}R)$ satisfies
  $\bigcap_{\ell=0}^sB(p_{i_\ell},10^{j-n-3}R)\subset
  B(p_{i_0},10^{j-n-2}R/2)\subset\bigcap_{\ell=0}^sB(p_{i_\ell},10^{j-n-2}R)$
  (triangle inequality), each summand is bounded by
  $\mathrm{cont}\,B(p_{i_0},10^{j-n-2}R/2)\le\mathcal{C}(k+1)$, since
  $B(p_{i_0},10^{j-n-2}R/2)\in\mathcal{B}(k+1)$ by the previous paragraph. The
  number of such intersections (indices $j\le n$, subsets of size
  $\le2^n\cdot10^{n(n+3)}$) is bounded by a constant $C(n)$ depending only on
  $n$, giving $\mathcal{C}(k)=\mathrm{cont}\,B(p,R)\le C(n)\mathcal{C}(k+1)$.
\end{proof}

\begin{theorem}[Thm. 12.5.1 — Gromov's finiteness theorem, 1978 and 1981]
  \label{thm:pet-ch12-betti-number-finiteness}
  \lean{PetersenLib.gromovFinitenessTheorem}
  \uses{thm:pet-ch12-toponogov,lem:pet-ch12-corank-bound,lem:pet-ch12-content-ratio-bound}
  There is $C(n)$ such that any complete $M$ with $\sec\ge0$ satisfies:
  \begin{enumerate}
    \item $\pi_1(M)$ can be generated by $\le C(n)$ elements;
    \item for every field $\mathbb{F}$,
      $\sum_{i=0}^n b_i(M,\mathbb{F})=\sum_{i=0}^n\dim H_i(M,\mathbb{F})\le C(n)$.
  \end{enumerate}
  \source{petersen:page-0483,petersen:page-0484,petersen:page-0492}
\end{theorem}
\begin{proof}
  \uses{thm:pet-ch12-betti-number-finiteness,thm:pet-ch12-toponogov,lem:pet-ch12-corank-bound,lem:pet-ch12-content-ratio-bound,lem:pet-ch12-gromov-critical-point-estimate}
  \emph{Part (1).} Let $\pi_1(M)$ act by deck transformations on the universal
  cover $\tilde M$, fix $p$, and inductively pick generators $g_1,g_2,\dots$
  with $|pg_k(p)|\le|pg(p)|$ for all $g\in\pi_1(M)-\{g_1,\dots,g_{k-1}\}$. If
  $k<l$ and $\angle(\overline{pg_k(p)},\overline{pg_l(p)})<\pi/3$, the hinge
  version of \cref{thm:pet-ch12-toponogov} gives
  \[
    |g_l(p)g_k(p)|^2<|pg_k(p)|^2+|pg_l(p)|^2-|pg_k(p)||pg_l(p)|\le|pg_l(p)|^2,
  \]
  so $|p(g_l^{-1}g_k)(p)|<|pg_l(p)|$, contradicting the choice of $g_l$. Hence
  the directions $\overline{pg_k(p)}$ pairwise make angle $\ge\pi/3$, so as in
  \cref{lem:pet-ch12-gromov-critical-point-estimate} there are at most
  $6^{n-1}$ of them, giving a bounded generating set.

  \emph{Part (2).} By \cref{lem:pet-ch12-corank-bound} the corank of any set is
  $\le100^n=:K$. Applying \cref{lem:pet-ch12-content-ratio-bound} repeatedly,
  $\mathrm{cont}\,M=\mathcal{C}(0)\le\mathcal{C}(k)\cdot C(n)^k$ for $k\le K$. If
  $\mathcal{B}(k)$ contains an incompressible ball, the proof of
  \cref{lem:pet-ch12-content-ratio-bound} produces a ball in $\mathcal{B}(k+1)$;
  taking $k=K$ (the largest corank actually attained in $M$), every ball in
  $\mathcal{B}(K)$ must be compressible, hence (compressing repeatedly) have
  minimal content $1$, so $\mathcal{C}(K)=1$. For $R$ large enough,
  $B(p,R)$ contains all the topology of $M$, so
  $\mathrm{cont}\,B(p,R)=\sum_ib_i(M,\mathbb{F})$, while by
  \cref{lem:pet-ch12-gromov-critical-point-estimate} the corank of $B(p,R)$ is
  $0$ for $R$ large. Hence
  $\sum_ib_i(M,\mathbb{F})=\mathrm{cont}\,M=\mathcal{C}(0)\le C(n)^K$.
\end{proof}

\section{Homotopy Finiteness}

\begin{lemma}[Lem. 12.6.2 — existence of a uniform regularity angle]
  \label{lem:pet-ch12-alpha-regular-existence}
  \lean{PetersenLib.uniformAlphaRegularityRadius}
  \uses{thm:pet-ch12-toponogov,def:pet-ch12-alpha-regular}
  For an integer $n>1$ and $v,D,k>0$, there are
  $\alpha=\alpha(n,D,v,k)\in(0,\pi/2)$ and $\delta=\delta(n,D,v,k)>0$ such that
  if $(M,g)$ is a Riemannian $n$-manifold with $\mathrm{diam}\le D$,
  $\mathrm{vol}\ge v$, $\sec\ge-k^2$, and $p,q\in M$ satisfy $|pq|\le\delta$,
  then $p$ is $\alpha$-regular for $q$ or $q$ is $\alpha$-regular for $p$.
  \source{petersen:page-0493,petersen:page-0494,petersen:page-0495,petersen:page-0496}
\end{lemma}
\begin{proof}
  \uses{lem:pet-ch12-alpha-regular-existence,thm:pet-ch12-toponogov}
  Argue by contradiction (a suggestion of Cheeger); take $k=-1$ for simplicity.
  Suppose $p,q$ are mutually not $\alpha$-regular with $|pq|\le\delta$. Then the
  sets $\overrightarrow{pq}\subset T_pM$, $\overrightarrow{qp}\subset T_qM$ of
  unit directions of segments from $p$ to $q$ (resp.\ $q$ to $p$) are
  $(\pi-\alpha)$-dense in the unit spheres. Since $t\mapsto\mathrm{vol}\,B(A,t)/v(n-1,1,t)$
  is nonincreasing for $A\subset S^{n-1}$, for a $(\pi-\alpha)$-dense
  $A\subset S^{n-1}$,
  \[
    \mathrm{vol}(S^{n-1}-B(A,\alpha))\le\mathrm{vol}\,S^{n-1}\cdot
    \frac{v(n-1,1,\pi-\alpha)-v(n-1,1,\alpha)}{v(n-1,1,\pi-\alpha)}.
  \]
  Choose $\alpha<\pi/2$ so small that this bound, multiplied by
  $\int_0^D\mathrm{sn}_k(t)^{n-1}\,dt$, equals $v/6$; then the two cones
  $B^{S^{n-1}-B(\overrightarrow{pq},\alpha)}(p,D)$ and
  $B^{S^{n-1}-B(\overrightarrow{qp},\alpha)}(q,D)$ (of directions avoiding the
  respective dense sets) each have volume $\le v/6$. Choose $r>0$ with
  $v(n,-1,r)=v/6$.

  We show that for $\delta$ small,
  \[
    M=B(p,r)\cup B(q,r)\cup B^{S^{n-1}-B(\overrightarrow{pq},\alpha)}(p,D)
    \cup B^{S^{n-1}-B(\overrightarrow{qp},\alpha)}(q,D),
  \]
  which forces $v\le\mathrm{vol}\,M\le4\cdot v/6<v$, a contradiction. If
  $x\notin B^{S^{n-1}-B(\overrightarrow{pq},\alpha)}(p,D)$, there is a hinge
  $\overline{xp}$, $\overline{pq}$ of angle $\le\alpha$ at $p$, so the hinge
  version of \cref{thm:pet-ch12-toponogov} (case $k=-1$) gives
  $\cosh|xq|\le\cosh|pq|\cosh|xp|-\sinh|pq|\sinh|xp|\cos\alpha$; if also
  $x\notin B^{S^{n-1}-B(\overrightarrow{qp},\alpha)}(q,D)$, symmetrically
  $\cosh|xp|\le\cosh|pq|\cosh|xq|-\sinh|pq|\sinh|xq|\cos\alpha$. If in addition
  $|xp|,|xq|>r$, adding $\cosh|xp|$ (resp.\ $\cosh|xq|$) to both estimates and
  using $|xp|,|xq|\le D$ gives
  \[
    \cosh|xq|\le\cosh|xp|+f(|pq|),\qquad\cosh|xp|\le\cosh|xq|+f(|pq|),
  \]
  where $f(t)=(\cosh t-1)\cosh D-\sinh t\sinh r\cos\alpha=(-\sinh
  r\cos\alpha)\,t+O(t^2)$. Since $f(0)=0$ and $f'(0)<0$, there is
  $\delta(D,r,\alpha)>0$ with $f(|pq|)<0$ for $|pq|\le\delta$, which makes both
  displayed inequalities contradictory ($\cosh|xq|<\cosh|xp|$ and
  $\cosh|xp|<\cosh|xq|$ simultaneously). Hence no such $x$ with $|xp|,|xq|>r$
  exists outside both cones, i.e.\ the four sets cover $M$, completing the
  contradiction and proving the lemma.
\end{proof}

\begin{corollary}[Cor. 12.6.3 — $\alpha$-regularity of the diagonal]
  \label{cor:pet-ch12-diagonal-alpha-regular}
  \lean{PetersenLib.diagonalAlphaRegular}
  \uses{lem:pet-ch12-alpha-regular-existence}
  With $\alpha,\delta$ as in \cref{lem:pet-ch12-alpha-regular-existence}, in
  $M\times M$ with the product metric let $\Delta=\{(x,x)\mid x\in M\}$. Any
  point within distance $\delta/\sqrt2$ of $\Delta$ is $\alpha$-regular for
  $\Delta$.
  \source{petersen:page-0496}
\end{corollary}
\begin{proof}
  \uses{cor:pet-ch12-diagonal-alpha-regular,lem:pet-ch12-alpha-regular-existence}
  Geodesics in $M\times M$ are pairs $(c_1,c_2)$ of geodesics in $M$; since
  $T_{(p,p)}\Delta=\{(v,v)\}$ has normal space
  $T^\perp_{(p,p)}\Delta=\{(v,-v)\}$, a segment $(c_1,c_2):[a,b]\to M\times M$
  from $(p,q)$ realizing $|(p,q)\Delta|$ satisfies
  $\dot c_1(b)=-\dot c_2(b)$, so $c_1,c_2$ join at $c_1(b)=c_2(b)$ into a
  geodesic from $p$ to $q$ in $M$, necessarily a segment (a shorter curve from
  $p$ to $q$ would bisect into a shorter curve from $(p,q)$ to $\Delta$). This
  gives a bijection between segments $\overline{pq}$ in $M$ and segments from
  $(p,q)$ to $\Delta$, with $\sqrt2\,|(p,q)\Delta|=|pq|$. Hence if
  $|(p,q)\Delta|\le\delta/\sqrt2$ then $|pq|\le\delta$, and by
  \cref{lem:pet-ch12-alpha-regular-existence} $p$ is $\alpha$-regular for $q$ or
  vice versa; translating this regularity through the above bijection makes
  $(p,q)$ $\alpha$-regular for $\Delta$.
\end{proof}

\begin{corollary}[Cor. 12.6.4 — nearby maps are homotopic]
  \label{cor:pet-ch12-close-maps-homotopic}
  \lean{PetersenLib.closeMapsAreHomotopic}
  \uses{cor:pet-ch12-diagonal-alpha-regular,lem:pet-ch12-retraction-lemma}
  With $\delta$ as in \cref{lem:pet-ch12-alpha-regular-existence}, if
  $f_0,f_1:X\to M$ are continuous with $|f_0(x)f_1(x)|<\delta$ for all $x\in X$,
  then $f_0,f_1$ are homotopic.
  \source{petersen:page-0496,petersen:page-0497}
\end{corollary}
\begin{proof}
  \uses{cor:pet-ch12-close-maps-homotopic,cor:pet-ch12-diagonal-alpha-regular,lem:pet-ch12-retraction-lemma}
  By \cref{cor:pet-ch12-diagonal-alpha-regular}, every point of $M\times M$
  within distance $\delta/\sqrt2$ of $\Delta$ is $\alpha$-regular for $\Delta$,
  so \cref{lem:pet-ch12-retraction-lemma} (with $K=\Delta$) provides a
  continuous deformation retraction of this neighborhood onto $\Delta$, moving
  each point a distance $\le\big(\sqrt2/\cos\alpha\big)|(p,q)\Delta|$; write
  $C=\sqrt2/\cos\alpha$. Translating through the bijection of
  \cref{cor:pet-ch12-diagonal-alpha-regular} between segments to $\Delta$ and
  segments in $M$, one obtains, for each pair $p,q\in M$ with $|pq|<\delta$, a
  curve $t\mapsto H(p,q,t)$, $0\le t\le1$, from $p$ to $q$ of length
  $\le C|pq|$, depending continuously on $(p,q,t)$. Given $f_0,f_1$ as
  hypothesized, $(x,t)\mapsto H(f_0(x),f_1(x),t)$ is then a homotopy from $f_0$
  to $f_1$.
\end{proof}

\begin{lemma}[Lem. 12.6.5 — extending nearby points to a simplex]
  \label{lem:pet-ch12-simplex-extension}
  \lean{PetersenLib.simplexExtensionLemma}
  \uses{cor:pet-ch12-close-maps-homotopic}
  With $C=\sqrt2/\cos\alpha$ and $\delta$ as above, suppose
  $p_0,\dots,p_k\in B(p,r)\subset M$ and $2r(C^k-1)/(C-1)<\delta$. Then there is
  a continuous map $f:\Delta^k\to B\big(p,r+2rC(C^k-1)/(C-1)\big)$ with
  $f(e_i)=p_i$, natural with respect to restriction to faces.
  \source{petersen:page-0497,petersen:page-0498,petersen:page-0499}
\end{lemma}
\begin{proof}
  \uses{lem:pet-ch12-simplex-extension,cor:pet-ch12-close-maps-homotopic}
  Induct on $k$; for $k=0$ there is nothing to show. Assume the map
  $f:\Delta^k\to B\big(p,2rC(C^k-1)/(C-1)+r\big)$ with $f(e_i)=p_i$,
  $i=0,\dots,k$, has been constructed for $k+1$ points $p_0,\dots,p_k$, and let
  $p_{k+1}\in B(p,r)$ be given. Define, using the map $H$ of
  \cref{cor:pet-ch12-close-maps-homotopic} (with the roles of $f_0,f_1$ played
  by the constant map at $p_{k+1}$ and by $f$ composed with radial projection
  from the new vertex),
  \[
    \tilde f(x^0,\dots,x^{k+1})=H\Big(f\Big(\tfrac{x^0}{\sum_{i\le k}x^i},\dots,\tfrac{x^k}{\sum_{i\le k}x^i}\Big),\,p_{k+1},\,x^{k+1}\Big).
  \]
  This is well defined and continuous provided
  $\big|p_{k+1}\,f(\cdots)\big|\le\big|p\,f(\cdots)\big|+|pp_{k+1}|\le
  \big(2rC\tfrac{C^k-1}{C-1}+r\big)+r=2r\tfrac{C^{k+1}-1}{C-1}<\delta$, which
  holds by hypothesis, so \cref{cor:pet-ch12-close-maps-homotopic} applies.
  Moreover $|p\tilde f(\cdot)|\le|pp_{k+1}|+|p_{k+1}\tilde f(\cdot)|\le
  r+2rC\tfrac{C^{k+1}-1}{C-1}$, completing the induction; restriction to any
  face spanned by a subset of the vertices agrees with the map constructed
  directly on that face, since the construction only uses the coordinates of
  the vertices involved.
\end{proof}

\begin{lemma}[Lem. 12.6.6 — homotopy equivalence from Gromov--Hausdorff closeness]
  \label{lem:pet-ch12-gh-close-homotopy-equivalent}
  \lean{PetersenLib.ghCloseImpliesHomotopyEquivalent}
  \uses{lem:pet-ch12-simplex-extension,cor:pet-ch12-close-maps-homotopic}
  There is $\varepsilon=\varepsilon(n,k,v,D)>0$ such that if $(M,g_1)$,
  $(N,g_2)$ are Riemannian $n$-manifolds with $\mathrm{diam}\le D$,
  $\mathrm{vol}\ge v$, $\sec\ge-k^2$, and Gromov--Hausdorff distance
  $d_{GH}(M,N)<\varepsilon$, then $M$ and $N$ are homotopy equivalent.
  \source{petersen:page-0499,petersen:page-0500}
\end{lemma}
\begin{proof}
  \uses{lem:pet-ch12-gh-close-homotopy-equivalent,lem:pet-ch12-simplex-extension,cor:pet-ch12-close-maps-homotopic}
  Realize $M,N$ as $\varepsilon$-Hausdorff-close subsets of a common metric
  space; $\varepsilon$ will be fixed through the construction. Triangulate both
  manifolds so every simplex lies in a ball of radius $\varepsilon$. Map the
  vertices $\{p_i\}\subset M$ of the triangulation to $\{q_i\}\subset N$ with
  $|p_iq_i|<\varepsilon$; a simplex $(p_{i_0},\dots,p_{i_n})$ lies in some
  $B(x,\varepsilon)$, so $\{q_{i_0},\dots,q_{i_n}\}\subset B(q_{i_0},4\varepsilon)$.
  If $8\varepsilon(C^n-1)/(C-1)<\delta$, \cref{lem:pet-ch12-simplex-extension}
  defines a map $f$ on each simplex, and these agree on common faces (by
  naturality of the construction), giving a continuous $f:M\to N$.

  For $x$ in the face spanned by $(p_{i_0},\dots,p_{i_n})$,
  \[
    |xf(x)|\le|xp_{i_0}|+|p_{i_0}f(x)|\le2\varepsilon+\varepsilon+|q_{i_0}f(x)|
    \le3\varepsilon+4\varepsilon+8\varepsilon C\tfrac{C^n-1}{C-1}
    =7\varepsilon+8\varepsilon C\tfrac{C^n-1}{C-1}.
  \]
  Construct $g:N\to M$ symmetrically with the same bound
  $|yg(y)|\le7\varepsilon+8\varepsilon C(C^n-1)/(C-1)$. Then
  \[
    |y(f\circ g)(y)|\le|yg(y)|+|g(y)(f\circ g)(y)|\le14\varepsilon+16\varepsilon
    C\tfrac{C^n-1}{C-1},
  \]
  and symmetrically for $|x(g\circ f)(x)|$. As long as
  $14\varepsilon+16\varepsilon C(C^n-1)/(C-1)<\delta$,
  \cref{cor:pet-ch12-close-maps-homotopic} shows $f\circ g\simeq\mathrm{id}_N$
  and $g\circ f\simeq\mathrm{id}_M$, so $f,g$ are homotopy inverses and
  $M\simeq N$. Since $C=\sqrt2/\cos\alpha$ and $\delta$ are given explicitly by
  \cref{lem:pet-ch12-alpha-regular-existence} in terms of $n,k,v,D$, this
  produces an explicit $\varepsilon(n,k,v,D)$.
\end{proof}

\begin{theorem}[Thm. — Grove--Petersen homotopy finiteness theorem, 1988]
  \label{thm:pet-ch12-homotopy-finiteness}
  \lean{PetersenLib.groveePetersenHomotopyFiniteness}
  \uses{lem:pet-ch12-gh-close-homotopy-equivalent}
  For an integer $n>1$ and $v,D,k>0$, the class of Riemannian $n$-manifolds with
  $\mathrm{diam}\le D$, $\mathrm{vol}\ge v$, $\sec\ge-k^2$ contains only finitely
  many homotopy types.
  \source{petersen:page-0493,petersen:page-0501}
\end{theorem}
\begin{proof}
  \uses{thm:pet-ch12-homotopy-finiteness,lem:pet-ch12-gh-close-homotopy-equivalent}
  The class in question is precompact in the Gromov--Hausdorff distance, since
  an upper diameter bound together with a lower sectional (hence Ricci)
  curvature bound gives, via Gromov's compactness criterion, a uniform bound on
  the number of $\varepsilon$-balls needed to cover any member of the class,
  for every $\varepsilon>0$. Covering this precompact class by finitely many
  Gromov--Hausdorff balls of radius $\varepsilon(n,k,v,D)$ as in
  \cref{lem:pet-ch12-gh-close-homotopy-equivalent}, every two members of the
  class lying in the same such ball are homotopy equivalent by that lemma.
  Hence the class contains at most as many homotopy types as the (finite)
  number of covering balls.
\end{proof}

\section{Further Study}

Several texts partially cover or expand the material of this chapter: the
surveys by Grove, by Abresch--Meyer, Colding, Greene, and Zhu, by Cheeger, and
by Karcher. The most notable omission from this chapter is the Abresch--Gromoll
theorem and other applications of the excess function; the surveys by Zhu and
by Cheeger cover this material.

\section{Exercises}

\begin{remark}[Exercise 12.8.1]
  \label{rem:pet-ch12-ex-12-8-1}
  \lean{PetersenLib.exercise_12_8_1}
  \uses{thm:pet-ch12-quarter-pinched-sphere,def:pet-ch12-alpha-regular}
  Let $(M,g)$ be a closed simply connected positively curved manifold. Show
  that if $M$ contains a totally geodesic closed hypersurface, then $M$ is
  homeomorphic to a sphere. (Hint: first show the hypersurface is orientable,
  then that the signed distance function to it has only two critical points, a
  maximum and a minimum; this also shows it suffices to assume
  $H^1(M,\mathbb{Z}_2)=0$.)
  \source{petersen:page-0501}
\end{remark}

\begin{remark}[Exercise 12.8.2]
  \label{rem:pet-ch12-ex-12-8-2}
  \lean{PetersenLib.exercise_12_8_2}
  \uses{thm:pet-ch12-soul-theorem}
  Let $(M,g)$ be complete noncompact with $\sec\ge0$ and soul $S\subset M$.
  \begin{enumerate}
    \item If $X\in TS$ and $V\in T^\perp S$, show $\sec(X,V)=0$.
    \item If $S$ has codimension $1$ with trivial normal bundle, show
      $(M,g)=(S\times\mathbb{R},g_S+dr^2)$.
    \item If $S$ has codimension $1$ with nontrivial normal bundle, show a
      double cover splits as in (2).
  \end{enumerate}
  \source{petersen:page-0502}
\end{remark}

\begin{remark}[Exercise 12.8.3]
  \label{rem:pet-ch12-ex-12-8-3}
  \lean{PetersenLib.exercise_12_8_3}
  \uses{thm:pet-ch12-toponogov}
  Show that the solution of $\ddot\xi=-(1+\varepsilon)\xi$,
  $\xi(0)=\psi(0)>0$, $\dot\xi(0)=\dot\psi(0)>0$ is
  \[
    \xi(t)=\sqrt{\psi(0)^2+\frac{\dot\psi(0)^2}{1+\varepsilon}}\,
    \sin\Big(\sqrt{1+\varepsilon}\,t+\arctan\Big(\frac{\psi(0)\sqrt{1+\varepsilon}}{\dot\psi(0)}\Big)\Big).
  \]
  \source{petersen:page-0502}
\end{remark}

\begin{remark}[Exercise 12.8.4]
  \label{rem:pet-ch12-ex-12-8-4}
  \lean{PetersenLib.exercise_12_8_4}
  \uses{thm:pet-ch12-toponogov}
  Show that the converse of Toponogov's theorem holds: if for some $k$ the
  conclusion of \cref{thm:pet-ch12-toponogov} holds when hinges (or triangles)
  are compared to the same objects in $S_k^2$, then $\sec\ge k$.
  \source{petersen:page-0502}
\end{remark}

\begin{remark}[Exercise 12.8.5]
  \label{rem:pet-ch12-ex-12-8-5}
  \lean{PetersenLib.exercise_12_8_5}
  \uses{thm:pet-ch12-toponogov}
  Let $(M,g)$ be complete. Show that if all sectional curvatures on $B(p,2R)$
  are $\ge k$, then Toponogov's comparison theorem
  (\cref{thm:pet-ch12-toponogov}) holds for hinges and triangles in $B(p,R)$.
  \source{petersen:page-0502}
\end{remark}

\begin{remark}[Exercise 12.8.6]
  \label{rem:pet-ch12-ex-12-8-6}
  \lean{PetersenLib.exercise_12_8_6}
  \uses{thm:pet-ch12-toponogov}
  Let $(M,g)$ be complete with $\sec\ge k$ and $f(x)=|xq|-|xp|$, with $|xq|$,
  $|xp|$ smooth at $x$. Show that \cref{thm:pet-ch12-toponogov} bounds
  $|\nabla f|_x$ from below in terms of the distance from $x$ to a segment from
  $q$ to $p$.
  \source{petersen:page-0502}
\end{remark}

\begin{remark}[Exercise 12.8.7 — Heintze and Karcher]
  \label{rem:pet-ch12-ex-12-8-7}
  \lean{PetersenLib.exercise_12_8_7}
  \uses{thm:pet-ch12-toponogov,lem:pet-ch12-alpha-regular-existence}
  Let $c\subset(M,g)$ be a geodesic in a Riemannian $n$-manifold with
  $\sec\ge-k^2$, and let $T(c,R)$ be the normal tube of radius $R$ around $c$
  (points joined to $c$ by a perpendicular segment of length $\le R$). In
  adapted tube coordinates $(r,s,\theta)$ the metric satisfies, as $r\to0$,
  \[
    g(r)=\begin{pmatrix}1&&\\&\ddots&\\&&0\end{pmatrix}
    +\begin{pmatrix}0&&\\&\ddots&\\&&1\end{pmatrix}r^2+O(r^3),
  \]
  with the first matrix having a single $0$ in the $s$-direction entry and the
  second a single $0$ there. Use the lower curvature bound to bound the volume
  density on the tube from above, and conclude
  $\mathrm{vol}\,T(c,R)\le f(n,k,R,L(c))$ for a continuous $f$ with
  $f\to0$ as $L(c)\to0$. Use this to prove Lem.\ 11.4.9 of the source and
  \cref{lem:pet-ch12-alpha-regular-existence}, showing Toponogov's theorem is
  not needed for the latter.
  \source{petersen:page-0502,petersen:page-0503}
\end{remark}

\begin{remark}[Exercise 12.8.8]
  \label{rem:pet-ch12-ex-12-8-8}
  \lean{PetersenLib.exercise_12_8_8}
  \uses{thm:pet-ch12-soul-theorem}
  Show that any vector bundle over $S^2$ admits a complete metric of
  nonnegative sectional curvature. (Hint: $k$-dimensional bundles over $S^2$
  are classified by $\pi_1(SO(k))$; there is one $1$-dimensional bundle, the
  $2$-dimensional bundles are parametrized by $\mathbb{Z}$, and for $k>2$ there
  are two $k$-dimensional bundles.)
  \source{petersen:page-0503}
\end{remark}

\begin{remark}[Exercise 12.8.9]
  \label{rem:pet-ch12-ex-12-8-9}
  \lean{PetersenLib.exercise_12_8_9}
  \uses{thm:pet-ch12-toponogov}
  Use Toponogov's theorem (\cref{thm:pet-ch12-toponogov}) to show that the
  Busemann function $b_r$ is convex when $\sec\ge0$.
  \source{petersen:page-0503}
\end{remark}
